\chapter{Thực nghiệm và Đánh giá hiệu năng}
\label{chap:evaluation}

Chương này trình bày quá trình thực nghiệm đo hiệu năng của Ordering Service ---
thành phần trọng tâm của hệ thống và cũng là nơi quyết định trần thông lượng đóng
block. Mục tiêu là (i)~xác định \textit{điểm bão hòa} của dịch vụ, (ii)~đo ảnh
hưởng của từng \textbf{tham số cấu hình} (kích thước block, chu kỳ đề xuất block)
và của từng \textbf{tham số tải} (tốc độ gửi, số luồng gửi), rồi (iii)~từ đó rút
ra bộ tham số cho hiệu năng tốt nhất. Toàn bộ số liệu trong chương được sinh ra
bằng công cụ \texttt{loadgen} đã mô tả ở Chương~\ref{chap:methodology}
(Mục~\ref{sec:measurement}), chạy trên cùng một máy để bảo đảm tính so sánh được.

\section{Mục tiêu và câu hỏi thực nghiệm}
\label{sec:eval-goals}

Ở Mục~\ref{sec:measurement}, khóa luận đã phân biệt ba lớp đo: lớp client (k6),
lớp Core (submit), và lớp ledger (commit). Chương này tập trung hoàn toàn vào
\textbf{lớp đóng block của Ordering Service} --- tức tốc độ mà leader lấy giao
dịch khỏi TxPool, tạo block, đạt đồng thuận và phát block đã commit qua deliver
stream. Việc cô lập đúng lớp này cho phép trả lời các câu hỏi sau mà không bị
nhiễu bởi WASM, ký endorsement hay độ trễ mirror PostgreSQL:

\begin{itemize}
  \item \textbf{Q1 --- Điểm bão hòa.} Khi tăng dần tốc độ gửi, tại đâu tốc độ
  commit ngừng bám theo tốc độ gửi? Trần thông lượng đóng block là bao nhiêu?
  \item \textbf{Q2 --- Kích thước block.} Tham số \texttt{auto\_propose\_block\_size}
  ảnh hưởng thế nào tới thông lượng (tx/s) và tần suất đóng block (blocks/s)?
  \item \textbf{Q3 --- Chu kỳ đề xuất block.} Tham số \texttt{auto\_propose\_interval}
  (block poll) có vai trò gì, và trong điều kiện tải nào nó thực sự chi phối?
  \item \textbf{Q4 --- Số luồng gửi.} Tham số \texttt{-workers} của loadgen ảnh
  hưởng ra sao tới thông lượng đo được?
  \item \textbf{Q5 --- Cấu hình tốt nhất.} Bộ tham số nào cho thông lượng cao và
  ổn định, và nó đánh đổi gì với độ trễ đóng block?
\end{itemize}

\section{Thiết lập thực nghiệm}
\label{sec:eval-setup}

\subsection{Môi trường và công cụ}
\label{sec:eval-env}

Toàn bộ phép đo chạy trên một máy đơn (single-host), loadgen và orderer giao tiếp
qua libp2p trên giao diện loopback. Cấu hình môi trường được liệt kê ở
Bảng~\ref{tab:eval-env}.

\begin{table}[H]
\centering
\caption{Môi trường thực nghiệm}
\label{tab:eval-env}
\begin{tabular}{@{}p{4.5cm}p{8cm}@{}}
\toprule
\textbf{Thành phần} & \textbf{Chi tiết} \\
\midrule
CPU & Intel Core Ultra 7 255H (16 nhân) \\
Bộ nhớ & 16\,GB RAM \\
Hệ điều hành & Windows 11 (64-bit) \\
Runtime & Go 1.25 \\
Mạng & libp2p trên loopback (không độ trễ mạng) \\
Công cụ đo & \texttt{loadgen} (giao thức endorsement + deliver stream) \\
\bottomrule
\end{tabular}
\end{table}

\texttt{loadgen} gửi các giao dịch hợp đồng giả lập (\texttt{bench\_ping}) trực
tiếp tới orderer qua đúng giao thức endorsement mà Core Service dùng, đồng thời
subscribe deliver stream để đếm số block và số giao dịch được commit. Chỉ số báo
cáo chính là \textbf{During load} --- số giao dịch/block được commit \textit{trong
đúng cửa sổ gửi tải}, phản ánh thông lượng bền vững của orderer; các block commit
muộn trong pha drain không được tính vào chỉ số này để tránh làm loãng con số.

\subsection{Mô hình một node và ý nghĩa của phép đo}
\label{sec:eval-singlenode}

Các phép đo dùng một cluster \textbf{một node} (leader tự đạt đa số --- self
majority), nên bước đồng thuận commit diễn ra tức thời. Lựa chọn này là có chủ
đích: nó \textit{cô lập trần của đường ống đóng block} (TxPool $\rightarrow$ đề
xuất $\rightarrow$ tính Merkle/hash $\rightarrow$ commit $\rightarrow$ deliver)
khỏi độ trễ vòng ACK giữa nhiều node. Vì hai tham số trọng tâm cần tinh chỉnh ---
kích thước block và chu kỳ đề xuất --- chi phối cơ chế gom lô (batching) bất kể số
node, đo trên một node cho tín hiệu sạch nhất về tác động của chúng. Trong cụm
nhiều node, thông lượng sẽ thấp hơn do cộng thêm một số hạng RTT đồng thuận gần
như cố định; con số một node vì vậy đóng vai \textbf{cận trên} của năng lực lớp
ordering (xem bàn luận ở Mục~\ref{sec:eval-threats}).

Do các phép đo throughput nhạy với nhiễu (GC, lập lịch của hệ điều hành, chia sẻ
CPU), mỗi cấu hình được chạy lặp lại \textbf{ba lần} và báo cáo theo giá trị
\textbf{trung vị} (median) để giảm ảnh hưởng của các lần chạy dị thường. Mỗi lần
chạy gồm 15\,s gửi tải và 8\,s drain.

\subsection{Ma trận tham số}
\label{sec:eval-matrix}

Bốn nhóm thí nghiệm (sweep) được thiết kế để tách biệt tác động của từng tham số
(Bảng~\ref{tab:eval-matrix}). Trong mỗi nhóm, chỉ một tham số thay đổi, các tham
số còn lại được giữ cố định.

\begin{table}[H]
\centering
\caption{Ma trận tham số thực nghiệm (mỗi cấu hình chạy 3 lần)}
\label{tab:eval-matrix}
\begin{tabular}{@{}p{3.2cm}p{4.3cm}p{4.8cm}@{}}
\toprule
\textbf{Nhóm} & \textbf{Tham số quét} & \textbf{Giữ cố định} \\
\midrule
Tải gửi (TPS) & tps $\in \{2, 5, 10, 15, 20\}$k, $\infty$ & block=1000, iv=100ms, w=32 \\
Kích thước block & block $\in \{100, 250, 500, 1000, 2000, 4000\}$ & tps=$\infty$, iv=100ms, w=32 \\
Chu kỳ đề xuất & iv $\in \{20, 50, 100, 200, 500\}$ms & block=1000, tps=5000, w=32 \\
Số luồng gửi & w $\in \{8, 16, 32, 48, 64\}$ & block=500, tps=$\infty$, iv=100ms \\
\bottomrule
\end{tabular}
\end{table}

Lưu ý: nhóm ``chu kỳ đề xuất'' được đo ở \textit{tải dưới bão hòa} (tps=5000) ---
đây là điều kiện duy nhất mà chu kỳ đề xuất thực sự chi phối nhịp đóng block (giải
thích ở Mục~\ref{sec:eval-interval}); và nhóm ``số luồng gửi'' được đo tại
\texttt{block=500} --- một điểm bão hòa \textit{ổn định} thay vì \texttt{block=1000}
(lý do ở Mục~\ref{sec:eval-blocksize}).

\section{Kết quả}
\label{sec:eval-results}

\subsection{Q1 --- Điểm bão hòa theo tốc độ gửi}
\label{sec:eval-tps}

Bảng~\ref{tab:eval-tps} và Hình~\ref{fig:eval-tps} trình bày kết quả khi tăng dần
tốc độ gửi ở cấu hình mặc định (block=1000, iv=100ms, 32 luồng). Trong mọi lần
chạy, số giao dịch gửi lỗi (\texttt{failed}) đều bằng 0.

\begin{table}[H]
\centering
\caption{Ảnh hưởng của tốc độ gửi (block=1000, iv=100ms, 32 luồng)}
\label{tab:eval-tps}
\begin{tabular}{@{}rrrrr@{}}
\toprule
\textbf{TPS đặt} & \textbf{Gửi (tx/s)} & \textbf{Commit (tx/s)} & \textbf{Blocks/s} & \textbf{Tx/block} \\
\midrule
2.000    & 1.987  & 1.986  & 9,87  & 201 \\
5.000    & 4.984  & 4.965  & 9,86  & 503 \\
10.000   & 9.933  & 9.864  & 9,87  & 997 \\
15.000   & 14.915 & 14.855 & 14,86 & 999 \\
20.000   & 19.866 & 19.857 & 19,87 & 1.000 \\
$\infty$ & 73.775 & 73.509 & 73,51 & 1.000 \\
\bottomrule
\end{tabular}
\end{table}

\begin{figure}[H]
\centering
\begin{tikzpicture}
\begin{axis}[
    width=0.82\textwidth, height=6.2cm,
    xlabel={Tốc độ gửi đặt (tx/s)}, ylabel={Commit (tx/s)},
    xmin=0, xmax=22000, ymin=0,
    legend pos=north west, grid=both, grid style={gray!20},
    scaled ticks=false, tick label style={/pgf/number format/1000 sep={}},
]
\addplot[dashed, gray, domain=0:21000] {x};
\addlegendentry{Lý tưởng (commit = gửi)}
\addplot[mark=*, thick, blue] coordinates {
    (2000,1986)(5000,4965)(10000,9864)(15000,14855)(20000,19857)};
\addlegendentry{Commit đo được}
\end{axis}
\end{tikzpicture}
\caption{Tốc độ commit bám sát tốc độ gửi tới ít nhất 20.000 tx/s; trần thực đo ở
tải không giới hạn là ${\sim}73.500$ tx/s.}
\label{fig:eval-tps}
\end{figure}

\textbf{Nhận xét.} Trong khoảng 2.000--20.000 tx/s, tốc độ commit gần như trùng
tốc độ gửi (sai lệch $<1\%$): orderer xử lý kịp, chưa bão hòa. Điểm đáng chú ý là
\textbf{tần suất đóng block} bị ghim ở ${\sim}9,87$ blocks/s khi tps $\le 10$k.
Nguyên nhân là cơ chế đề xuất theo sự kiện (event-driven) mô tả ở
Chương~\ref{chap:methodology}: khi TxPool chưa gom đủ một lô đầy (1.000 tx) trong
một chu kỳ, leader chờ hết chu kỳ 100ms rồi mới flush --- cho đúng $1/0{,}1 = 10$
block mỗi giây. Vì thế số tx mỗi block tăng tuyến tính theo tải (201 $\to$ 503
$\to$ 997) trong khi blocks/s giữ nguyên. Khi tải vượt ngưỡng lấp đầy lô
(${\ge}15$k), pool đầy trước 100ms nên leader đề xuất ngay, blocks/s bắt đầu tăng
(14,86 rồi 19,87). Ở tải không giới hạn, hệ thống đạt trần
\textbf{${\sim}73.500$ tx/s} tương ứng ${\sim}73,5$ block/s --- đây là năng lực
đóng block thuần của một leader trên môi trường thử nghiệm.

\subsection{Q2 --- Ảnh hưởng của kích thước block}
\label{sec:eval-blocksize}

Bảng~\ref{tab:eval-blocksize} và Hình~\ref{fig:eval-blocksize} đo thông lượng ở
tải bão hòa (tps không giới hạn) khi thay đổi kích thước block.

\begin{table}[H]
\centering
\caption{Ảnh hưởng của kích thước block (tps=$\infty$, iv=100ms, 32 luồng)}
\label{tab:eval-blocksize}
\begin{tabular}{@{}rrrr@{}}
\toprule
\textbf{Block size} & \textbf{Commit (tx/s)} & \textbf{Blocks/s} & \textbf{Tx/block} \\
\midrule
100  & 74.546 & 745,5 & 100 \\
250  & \textbf{82.497} & 330,0 & 250 \\
500  & 77.774 & 155,6 & 500 \\
1.000 & 57.439 & 57,4 & 1.000 \\
2.000 & 67.307 & 33,7 & 2.000 \\
4.000 & 73.972 & 18,5 & 3.999 \\
\bottomrule
\end{tabular}
\end{table}

\begin{figure}[H]
\centering
\begin{tikzpicture}
\begin{axis}[
    width=0.82\textwidth, height=6.2cm,
    xlabel={Kích thước block (tx)}, ylabel={Commit (tx/s)},
    xmode=log, log basis x=2,
    xtick={100,250,500,1000,2000,4000},
    xticklabels={100,250,500,1000,2000,4000},
    ymin=40000, ymax=90000,
    grid=both, grid style={gray!20},
    scaled ticks=false, tick label style={/pgf/number format/1000 sep={}},
    x tick label style={rotate=45, anchor=east},
]
\addplot[mark=*, thick, red] coordinates {
    (100,74546)(250,82497)(500,77774)(1000,57439)(2000,67307)(4000,73972)};
\end{axis}
\end{tikzpicture}
\caption{Thông lượng commit theo kích thước block: đạt đỉnh ${\sim}82.500$ tx/s ở
block=250, với một vùng trũng lặp lại được ở block=1.000.}
\label{fig:eval-blocksize}
\end{figure}

\textbf{Nhận xét.} Tần suất đóng block tỉ lệ nghịch với kích thước block (blocks/s
$\times$ block size $\approx$ tx/s), đúng như kỳ vọng. Về thông lượng, tx/s giữ ở
mức cao (${\sim}67$--$82$k) trên phần lớn dải, \textbf{đạt đỉnh ${\sim}82.500$ tx/s
tại block=250}. Đáng chú ý, \texttt{block=1.000} --- giá trị mặc định --- lại là
một \textbf{vùng trũng lặp lại được} (${\sim}57.400$ tx/s, thấp hơn hai lân cận),
và các lần chạy tại điểm này cũng \textit{dao động mạnh nhất} (36k--83k giữa các
lần đo). Điều này gợi ý \texttt{block=1.000} nằm gần một \textit{ranh giới bất ổn}:
tốc độ lấp đầy lô dao động quanh đúng ngưỡng batch nên vòng đề xuất lúc thì chờ
hết chu kỳ, lúc thì flush ngay --- tạo hành vi hai chế độ (bimodal). Các kích thước
nhỏ hơn (250--500) hoặc lớn hơn (4.000) đều ổn định và cho thông lượng cao hơn.
Đây là một kết quả thực nghiệm quan trọng: \textbf{giá trị mặc định 1.000 vừa
không tối ưu về thông lượng vừa kém ổn định}.

\subsection{Q3 --- Vai trò của chu kỳ đề xuất block}
\label{sec:eval-interval}

Cần phân biệt hai chế độ. Ở \textit{tải bão hòa}, TxPool luôn có sẵn hơn một lô
đầy, nên vòng đề xuất theo sự kiện bỏ qua thời gian chờ chu kỳ --- do đó chu kỳ đề
xuất \textbf{gần như không tác động} tới thông lượng (thí nghiệm đo tại
block=1.000, tps=$\infty$ cho các giá trị dao động 56k--76k không theo xu hướng rõ
ràng, đúng với dự đoán rằng chu kỳ bị bỏ qua). Vai trò thực của chu kỳ chỉ lộ ra ở
\textit{tải dưới bão hòa}, khi pool chưa gom đủ lô: khi đó chu kỳ quyết định nhịp
đóng block và độ mịn gom lô. Bảng~\ref{tab:eval-interval} đo ở tps=5.000 (dưới
bão hòa).

\begin{table}[H]
\centering
\caption{Ảnh hưởng của chu kỳ đề xuất ở tải dưới bão hòa (tps=5.000, block=1.000, 32 luồng)}
\label{tab:eval-interval}
\begin{tabular}{@{}rrrr@{}}
\toprule
\textbf{Chu kỳ (ms)} & \textbf{Commit (tx/s)} & \textbf{Blocks/s} & \textbf{Tx/block} \\
\midrule
20  & 4.966 & 15,19 & 327 \\
50  & 4.966 & 11,87 & 418 \\
100 & 4.964 & 9,80  & 507 \\
200 & 4.932 & 4,93  & 1.000 \\
500 & 4.800 & 4,80  & 1.000 \\
\bottomrule
\end{tabular}
\end{table}

\begin{figure}[H]
\centering
\begin{tikzpicture}
\begin{axis}[
    width=0.82\textwidth, height=6cm,
    xlabel={Chu kỳ đề xuất (ms)}, ylabel={Giá trị},
    xmin=0, xmax=520, ymin=0,
    legend pos=north west, grid=both, grid style={gray!20},
    scaled ticks=false, tick label style={/pgf/number format/1000 sep={}},
]
\addplot[mark=*, thick, blue] coordinates {(20,327)(50,418)(100,507)(200,1000)(500,1000)};
\addlegendentry{Tx/block}
\addplot[mark=square*, thick, orange] coordinates {(20,15.19)(50,11.87)(100,9.80)(200,4.93)(500,4.80)};
\addlegendentry{Blocks/s}
\end{axis}
\end{tikzpicture}
\caption{Ở tải 5.000 tx/s: chu kỳ ngắn cho nhiều block nhỏ, chu kỳ dài cho ít
block lớn (bão hòa ở kích thước lô 1.000); thông lượng commit không đổi ${\sim}5.000$ tx/s.}
\label{fig:eval-interval}
\end{figure}

\textbf{Nhận xét.} Ở tải 5.000 tx/s, tốc độ commit giữ nguyên ${\sim}5.000$ tx/s
với mọi chu kỳ (orderer luôn theo kịp), nhưng \textbf{tần suất đóng block và số
tx/block thay đổi rõ rệt}: chu kỳ ngắn cho nhiều block nhỏ (độ trễ đóng block
thấp), chu kỳ dài cho ít block lớn (độ trễ cao hơn nhưng ít overhead trên mỗi
giao dịch). Nói cách khác, chu kỳ đề xuất là \textbf{núm điều chỉnh đánh đổi độ
trễ $\leftrightarrow$ độ mịn gom lô}, không phải núm điều chỉnh thông lượng.

\subsection{Q4 --- Ảnh hưởng của số luồng gửi}
\label{sec:eval-workers}

Để tách tác động của số luồng gửi khỏi sự bất ổn tại block=1.000, thí nghiệm này
đo tại điểm bão hòa ổn định \texttt{block=500}. Kết quả ở
Bảng~\ref{tab:eval-workers} và Hình~\ref{fig:eval-workers}.

\begin{table}[H]
\centering
\caption{Ảnh hưởng của số luồng gửi (block=500, tps=$\infty$, iv=100ms)}
\label{tab:eval-workers}
\begin{tabular}{@{}rrr@{}}
\toprule
\textbf{Số luồng} & \textbf{Commit (tx/s)} & \textbf{Blocks/s} \\
\midrule
8  & 80.196 & 160,4 \\
16 & 81.590 & 163,2 \\
32 & \textbf{82.780} & 165,6 \\
48 & 63.963 & 127,9 \\
64 & 65.984 & 132,0 \\
\bottomrule
\end{tabular}
\end{table}

\begin{figure}[H]
\centering
\begin{tikzpicture}
\begin{axis}[
    width=0.82\textwidth, height=6cm,
    xlabel={Số luồng gửi}, ylabel={Commit (tx/s)},
    xmin=0, xmax=68, ymin=50000, ymax=90000,
    xtick={8,16,32,48,64},
    grid=both, grid style={gray!20},
    scaled ticks=false, tick label style={/pgf/number format/1000 sep={}},
]
\addplot[mark=*, thick, teal] coordinates {(8,80196)(16,81590)(32,82780)(48,63963)(64,65984)};
\end{axis}
\end{tikzpicture}
\caption{Thông lượng theo số luồng gửi: vùng phẳng tối ưu ở 8--32 luồng (đỉnh tại
32, ${\sim}82.800$ tx/s), suy giảm khi vượt 32 do tranh chấp stream.}
\label{fig:eval-workers}
\end{figure}

\textbf{Nhận xét.} Số luồng gửi ảnh hưởng tới \textit{tốc độ nạp} (ingest) chứ
không phải năng lực đóng block của orderer; tuy nhiên nếu quá ít luồng thì không
bơm đủ tải để orderer bão hòa, còn quá nhiều luồng lại gây tranh chấp (mỗi luồng
giữ một stream libp2p) và làm giảm thông lượng đo được. Kết quả cho thấy một vùng
tối ưu ở giữa; ngoài vùng này thông lượng suy giảm.

\section{Cấu hình khuyến nghị}
\label{sec:eval-best}

Tổng hợp bốn nhóm thí nghiệm, bộ tham số cho hiệu năng cao và ổn định được đề
xuất ở Bảng~\ref{tab:eval-best}.

\begin{table}[H]
\centering
\caption{Bộ tham số khuyến nghị cho Ordering Service}
\label{tab:eval-best}
\begin{tabular}{@{}p{4.2cm}p{3cm}p{5.3cm}@{}}
\toprule
\textbf{Tham số} & \textbf{Khuyến nghị} & \textbf{Lý do} \\
\midrule
\texttt{auto\_propose\_block\_size} & 250--500 & Thông lượng đỉnh (${\sim}82.500$ tx/s ở 250), ổn định; tránh vùng trũng bất ổn ở 1.000 \\
\texttt{auto\_propose\_interval} & 100ms & Cân bằng độ trễ/gom lô ở tải thấp; không ảnh hưởng khi bão hòa \\
Số luồng gửi (loadgen) & 32 & Đủ bơm bão hòa mà chưa gây tranh chấp stream \\
\bottomrule
\end{tabular}
\end{table}

So với mặc định (block=1.000), chuyển sang \texttt{block=250}--\texttt{500} vừa
\textbf{tăng thông lượng} (${\sim}57$k $\to$ ${\sim}78$--$82$k tx/s, cải thiện
${\sim}35$--$44\%$) vừa \textbf{loại bỏ dao động} của điểm vận hành. Đánh đổi duy
nhất là số block nhiều hơn (blocks/s cao hơn), làm tăng chút overhead ở các thành
phần hạ nguồn (Committing Peer, mirror PostgreSQL); tuy nhiên các thành phần này
được thiết kế bất đồng bộ nên không nằm trên đường tới hạn của việc đóng block.

\section{Bàn luận và giới hạn}
\label{sec:eval-threats}

\textbf{Ordering không phải là nút thắt của hệ thống.} Trần đóng block thuần
(${\sim}73$--$82$k tx/s) cao hơn nhiều lần so với thông lượng end-to-end tham khảo
qua toàn pipeline (${\sim}3.900$ commit/s, Mục~\ref{sec:measurement}). Điều này
khẳng định lớp ordering \textit{không} phải cổ chai: giới hạn thực tế nằm ở
đường Core (WASM + ký endorsement) và ở ràng buộc ``một block in-flight'' của vòng
đồng thuận nhiều node (Mục~\ref{sec:opts}), chứ không ở năng lực gom lô/đóng block
của leader.

\textbf{Giới hạn của phép đo.} Các con số trên là \textbf{cận trên} của lớp
ordering, cần diễn giải trong phạm vi sau: (i)~cluster một node nên không có độ
trễ ACK đồng thuận giữa nhiều node; (ii)~loopback nên không có độ trễ mạng thực;
(iii)~trạng thái in-memory nên không có chi phí fsync xuống đĩa; (iv)~giao dịch
giả lập không thực thi WASM. Trong triển khai nhiều node trên mạng thật, thông
lượng sẽ thấp hơn do cộng thêm RTT đồng thuận và chi phí I/O; song thứ hạng tương
đối giữa các cấu hình (đặc biệt vùng trũng tại block=1.000 và đỉnh tại
block=250--500) được kỳ vọng giữ nguyên vì chúng bắt nguồn từ cơ chế gom lô độc
lập với số node. (v)~Phép đo thực hiện trên một máy chia sẻ tài nguyên nên còn
nhiễu; việc lặp ba lần và lấy trung vị nhằm giảm thiểu ảnh hưởng này.
